% Compile here: latexmk -pdf report.tex
\documentclass[12pt,a4paper]{article}
\usepackage{cmap}
\usepackage[utf8]{inputenc}
\usepackage[T2A]{fontenc}
\usepackage[english,russian]{babel}
\usepackage{amsmath,amssymb,bm}
\usepackage{graphicx,float,booktabs,geometry,array,needspace}
\usepackage[justification=raggedright,singlelinecheck=false]{caption}
\usepackage[expansion=false]{microtype}
\usepackage{xcolor}
\usepackage{hyperref}
\geometry{margin=2.2cm}
\setlength{\emergencystretch}{2em}
\AddToHook{cmd/subsection/before}{\Needspace{8\baselineskip}}
\hypersetup{colorlinks=true,linkcolor=black,citecolor=black,urlcolor=blue,
 pdftitle={Резервуарная и спиновая динамика поляритонного конденсата},pdfauthor={Коростелев Никита}}
\graphicspath{{../results/figures/report/}}
\makeatletter
\setlength{\@fptop}{0pt}
\setlength{\@fpsep}{12pt}
\setlength{\@fpbot}{0pt plus 1fil}
\makeatother
\newcommand{\reportfigure}[3]{%
\begin{figure}[H]
\centering
\IfFileExists{../results/figures/report/#1}{%
 \includegraphics[width=0.96\textwidth]{#1}
 \caption{#2}
}{%
 \fbox{\parbox[c][48mm][c]{0.9\textwidth}{\centering
 \textbf{Пространственный расчёт выполняется}\\[1em]
 Карты будут добавлены после завершения новой серии.\\
 Промежуточные поля не выдаются за установившийся результат.}}
 \caption{#2 \textbf{В этой предварительной версии карты ещё не включены: расчёт выполняется.}}
}
\label{#3}
\end{figure}}
\title{Резервуарная и спиновая динамика поляритонного конденсата}
\author{Коростелев Никита}
\date{}
\begin{document}
\maketitle
\clearpage

\section*{Цель работы}
\textit{Дополнение к отчёту о конденсации экситонных поляритонов.}
% REPORT_STATUS_START

% REPORT_STATUS_END

В предыдущем отчёте «Динамика и пространственная структура экситон-поля\-ри\-тонных конденсатов при нерезонансной накачке» рассматривался скалярный поляритонный конденсат, связанный с одним резервуаром. Здесь модель последовательно расширяется: вводится неактивный резервуар, затем две круговые компоненты конденсата и когерентная связь между ними.

Рассматриваются три вопроса: как дополнительный резервуар изменяет переходный процесс; как поляризация накачки влияет на плотности спиновых компонент; при каких значениях связи плотности и разность фаз становятся постоянными во времени. Основные параметры прежней модели сохраняются. Логика выбора новых коэффициентов приведена в приложении~\hyperref[sec:A]{A}.

\section{Активный и неактивный резервуары}\label{sec:1}

\subsection{Расширение скалярной модели}\label{sec:1.1}

В прежней модели поле конденсата \(\psi\) и плотность резервуара \(n_R\) удовлетворяли уравнениям

\begin{equation}\label{eq:1}
\begin{gathered}
i\hbar\frac{\partial\psi}{\partial t}
=
\left[
-K\nabla^2+g_C|\psi|^2
+g_R(n_R+\eta P_A)
+\frac{i\hbar}{2}(Rn_R-\gamma_C)
\right]\psi,\\[0.5em]
\frac{\partial n_R}{\partial t}
=
P_A-(\gamma_R+R|\psi|^2)n_R,
\qquad
K=\frac{\hbar^2}{2m}.
\end{gathered}
\end{equation}

Здесь \(P_A\) — скорость поступления частиц в резервуар. Обозначение отличается от прежнего \(P\), чтобы далее различать внешнюю накачку и питание активного резервуара.

Слагаемое \(g_R\eta P_A\) задаёт дополнительный энергетический сдвиг, который мгновенно следует за изменением накачки. Теперь этот вклад описывается через отдельную динамическую плотность \(n_I\). Неактивный резервуар получает частицы от внешней накачки \(P_{\mathrm{in}}\) и передаёт их активному резервуару \(n_A\) со скоростью \(Wn_I\). Активный резервуар непосредственно обеспечивает стимулированное поступление частиц в конденсат:

\begin{equation}\label{eq:2}
\begin{gathered}
i\hbar\frac{\partial\psi}{\partial t}
=
\left[
-K\nabla^2+g_C|\psi|^2+g_A n_A+g_I n_I
+\frac{i\hbar}{2}(Rn_A-\gamma_C)
\right]\psi,\\[0.5em]
\frac{\partial n_A}{\partial t}
=
Wn_I-(\gamma_A+R|\psi|^2)n_A,\\[0.5em]
\frac{\partial n_I}{\partial t}
=
P_{\mathrm{in}}-(\gamma_I+W)n_I.
\end{gathered}
\end{equation}

Оба резервуара сдвигают энергию конденсата. Прежнее слагаемое \(g_R\eta P_A\) в уравнение \eqref{eq:2} повторно не включается. Исходная скалярная модель и двухрезервуарное расширение опираются на работы \cite{ref1,ref2}.

Для согласования моделей выбираются

\begin{equation}\label{eq:3}
g_A=g_R,\qquad
\gamma_A=\gamma_R,\qquad
P_{\mathrm{in}}=\frac{\gamma_I+W}{W}P_A,
\qquad
\frac{g_I}{W}=g_R\eta.
\end{equation}

При постоянной накачке \(n_I^*=P_A/W\). Поэтому питание активного резервуара и дополнительный энергетический сдвиг совпадают с прежними значениями:

\[
Wn_I^*=P_A,\qquad g_I n_I^*=g_R\eta P_A.
\]

Таким образом, сохраняются стационарные решения скалярной модели. В пространственно однородном случае сохраняется и порог

\[
P_{\mathrm{th}}=\frac{\gamma_C\gamma_A}{R}
=0.43333\ \mu\mathrm m^{-2}\,\mathrm{ps}^{-1}.
\]

Это согласование не требует совпадения переходных процессов.

\subsection{Отклик на включение и выключение накачки}\label{sec:1.2}

Введём поток частиц в активный резервуар

\[
Q_A=Wn_I.
\]

Из уравнений \eqref{eq:2} и условий \eqref{eq:3} следует

\[
\tau_I\frac{dQ_A}{dt}=P_A-Q_A,
\qquad
\tau_I=\frac{1}{\gamma_I+W}.
\]

Вместо мгновенного изменения питания активного резервуара возникает отклик с конечным временем \(\tau_I\). Для основного набора параметров \(\tau_I=9.63\) пс.

Сравнение проводится при \(P_A=1.5P_{\mathrm{th}}\), пустых начальных резервуарах и одинаковой слабой затравке конденсата. Обе модели выходят на плотность

\[
n_C^*=\frac{P_A-P_{\mathrm{th}}}{\gamma_C}
=1.19048\ \mu\mathrm m^{-2}.
\]

Однако рост конденсата в двухрезервуарной модели начинается позже. Первое достижение половины стационарной плотности происходит при \(159.1\) пс в прежней модели и при \(188.0\) пс в новой. Эта задержка определяется всем переходным процессом и не равна непосредственно \(\tau_I\).

После выключения накачки неактивный резервуар некоторое время продолжает питать активный. В прежней модели внешнее питание активного резервуара прекращается сразу, но уже накопленные в нём частицы также могут поступать в конденсат.

\reportfigure{N1_reservoir_switching.png}{Однородная плотность конденсата \(n_C=|\psi|^2\) в моделях с одним и двумя резервуарами. \textbf{(a)} Рост после включения накачки \(P_A=1.5P_{\mathrm{th}}\); горизонтальный пунктир — общая стационарная плотность \(n_C^*=1.19048\ \mu\mathrm m^{-2}\). \textbf{(b)} Затухание после выключения при \(t_{\mathrm{off}}=1000\) пс; время отсчитывается от выключения. В обоих расчётах \(n_C(0)=10^{-4}\ \mu\mathrm m^{-2}\), резервуары первоначально пусты. Параметры приведены в приложении~\hyperref[sec:A]{A}.}{fig:1}



\clearpage
\section{Две спиновые компоненты конденсата}\label{sec:2}

\subsection{Уравнения модели}\label{sec:2.1}

В круговом базисе конденсат описывается полями \(\psi_+\) и \(\psi_-\), соответствующими двум круговым поляризациям. Их плотности равны \(n_\pm=|\psi_\pm|^2\). Для каждой компоненты вводятся активный и неактивный резервуары.

Используется система

\begin{equation}\label{eq:4}
\begin{gathered}
i\hbar\frac{\partial\psi_\sigma}{\partial t}
=
\left[
-K\nabla^2+V_\sigma
+\frac{i\hbar}{2}(Rn_{A\sigma}-\gamma_C)
\right]\psi_\sigma
-J\psi_{-\sigma},\\[0.5em]
V_\sigma
=
\alpha_1 n_\sigma+\alpha_2 n_{-\sigma}
+g_A n_{A\sigma}+g_I n_{I\sigma},\\[0.5em]
\frac{\partial n_{A\sigma}}{\partial t}
=
Wn_{I\sigma}
-(\gamma_A+Rn_\sigma)n_{A\sigma}
+\Gamma_{s,A}(n_{A,-\sigma}-n_{A\sigma}),\\[0.5em]
\frac{\partial n_{I\sigma}}{\partial t}
=
P_{\mathrm{in},\sigma}
-(\gamma_I+W)n_{I\sigma}
+\Gamma_{s,I}(n_{I,-\sigma}-n_{I\sigma}),
\qquad \sigma=\pm.
\end{gathered}
\end{equation}

Коэффициенты \(\alpha_1\) и \(\alpha_2\) задают энергетические сдвиги от поляритонов той же и противоположной поляризации. Слагаемые с \(\Gamma_s\) перераспределяют частицы между спиновыми компонентами резервуаров, сохраняя их суммарную плотность. Такая структура спиновой модели используется в работе \cite{ref3}.

В этой главе \(J=0\); когерентная связь рассматривается в разделе~\ref{sec:3}. Сдвиг от резервуара учитывается только для соответствующей круговой компоненты: перекрёстные резервуарные коэффициенты приняты равными нулю.

Все взаимодействия в \eqref{eq:4} локальны. Связь различных пространственных точек возникает через кинетический оператор \(-K\nabla^2\). Диффузия резервуаров в данной модели отсутствует.

\subsection{Зависимость плотностей от поляризации накачки}\label{sec:2.2}

Для однородной задачи зададим

\begin{equation}\label{eq:5}
P_{A,\pm}=sP_{\mathrm{th}}(1\pm\varepsilon),
\qquad
P_{\mathrm{in},\pm}
=
\frac{\gamma_I+W}{W}P_{A,\pm}.
\end{equation}

Здесь \(P_{A,\pm}\) обозначает перенормированный параметр внешней накачки. Фактический поток частиц в каждый активный резервуар равен \(Q_{A,\pm}=Wn_{I,\pm}\) и в общем случае отличается от \(P_{A,\pm}\) из-за спиновой релаксации в неактивном резервуаре.

Параметр \(\varepsilon\) характеризует дисбаланс накачки двух компонент. При его изменении суммарная накачка \(P_{A,+}+P_{A,-}=2sP_{\mathrm{th}}\) остаётся постоянной. Далее используется \(s=2\), то есть суммарная эффективная накачка равна \(4P_{\mathrm{th}}\).

При \(J=0\) и ненулевых стационарных плотностях обеих компонент условие баланса усиления и потерь даёт

\[
n_{A,+}^*=n_{A,-}^*=\frac{\gamma_C}{R}.
\]

Обозначив

\[
D_I=
\frac{\gamma_I+W}
{\gamma_I+W+2\Gamma_{s,I}},
\]

\noindent\begin{minipage}{\textwidth}
получаем

\begin{equation}\label{eq:6}
\begin{gathered}
n_{I,\pm}^*
=
\frac{sP_{\mathrm{th}}}{W}(1\pm\varepsilon D_I),\\[0.5em]
n_\pm^*
=
\frac{\gamma_A}{R}
\left[s(1\pm\varepsilon D_I)-1\right].
\end{gathered}
\end{equation}
\end{minipage}
\par\medskip

Эти выражения применимы, когда обе полученные плотности положительны:

\[
s(1-|\varepsilon|D_I)>1.
\]

При выбранных параметрах \(D_I=0.34180\), и условие выполняется для всего рассматриваемого диапазона \(-1\leq\varepsilon\leq1\).

Дисбаланс плотностей конденсата равен

\begin{equation}\label{eq:7}
z=
\frac{n_+-n_-}{n_++n_-}
=
\frac{s\varepsilon D_I}{s-1}.
\end{equation}

При \(s=2\) получаем \(z=0.68360\varepsilon\). Спиновое перераспределение в неактивном резервуаре уменьшает различие между потоками, питающими две компоненты. При этом активные резервуары имеют одинаковые стационарные плотности, а различие потоков проявляется в разных плотностях конденсата.

Численное интегрирование воспроизводит зависимости \eqref{eq:6}. Например, при \(\varepsilon=0.3\)

\[
n_+^*=2.86924,\qquad
n_-^*=1.89267\ \mu\mathrm m^{-2},
\qquad z=0.20508.
\]

\reportfigure{N2_circular_populations.png}{Стационарные плотности круговых компонент \(n_+\) и \(n_-\) в зависимости от дисбаланса накачки \(\varepsilon\). Суммарная эффективная накачка фиксирована: \(P_{A,+}+P_{A,-}=4P_{\mathrm{th}}\); когерентная связь отсутствует, \(J=0\). Линии показывают аналитические зависимости \eqref{eq:6}, точки — результаты независимого интегрирования для 21 значения \(\varepsilon\). Численные значения определены на интервале \(1500\)–\(2000\) пс. Для всех запусков используются одинаковые начальные плотности компонент и пустые резервуары.}{fig:2}



\subsection{Пространственные распределения при кольцевой накачке}\label{sec:2.3}

Для перехода к пространственной задаче используется кольцевая накачка

\begin{equation}\label{eq:8}
P_{A,\pm}(r)
=
s_{\mathrm{ring}}P_{\mathrm{th}}(1\pm\varepsilon)
\exp\!\left[-\frac{(r-r_0)^2}{d^2}\right],
\end{equation}

где \(r_0=10\) мкм, \(d=2.5\) мкм и \(s_{\mathrm{ring}}=3.16875\). Сравниваются \(\varepsilon=0\) и \(0.3\) при \(J=0\).

При равной накачке её амплитуда для каждой компоненты совпадает с амплитудой скалярного кольцевого расчёта из прошлого отчёта. Суммарная пиковая накачка здесь равна \(6.3375P_{\mathrm{th}}\) и одинакова в обоих спиновых расчётах.

Начальные поля выбираются как

\[
\psi_+(\mathbf r,0)
=
\psi_-(\mathbf r,0)
=
\frac{\psi_{\mathrm{seed}}(\mathbf r)}{\sqrt{2}},
\]

где \(\psi_{\mathrm{seed}}\) — прежняя пространственная затравка из 24 плоских волн с огибающей в области накачки. Все резервуары первоначально пусты.

Действие дисбаланса накачки оценивается по картам на общей абсолютной цветовой шкале и интегральным величинам

\[
N_\pm=\int n_\pm(\mathbf r)\,d^2r,
\qquad
Z=\frac{N_+-N_-}{N_++N_-}.
\]

Однородную формулу \eqref{eq:7} нельзя непосредственно переносить на эту задачу: пространственный перенос и неоднородность резервуаров учитываются явно.

% N3_RESULT_START
Карты на рис.~\ref{fig:3} соответствуют общему времени $t=2300$ пс. Интегральные величины вычислены по всей области $256\times256$ мкм.
\begin{center}\begin{tabular}{@{}crrrc@{}}\toprule
$\varepsilon$ & $N_+$ & $N_-$ & $Z$ & Установление\\\midrule
0 & 305.8816 & 305.8816 & 0.00000 & да\\
0.3 & 440.2656 & 169.6108 & 0.44379 & да\\
\bottomrule\end{tabular}\end{center}
Обе компоненты и соответствующие резервуары прошли проверку установления на двух последовательных окнах по 100 пс.
При симметричной накачке и одинаковых начальных полях плотности компонент совпадают. Это контроль симметричной подготовки, а не доказательство выбора равных плотностей из любых начальных условий.
При $\varepsilon=0.3$ дисбаланс накачки прежде всего меняет населённости компонент: $Z=0.44379$. Пространственные профили остаются близкими, но не совпадают. Их относительное различие после нормировки каждого на соответствующую населённость определяется как
\[\begin{gathered}\widehat n_\sigma=\frac{n_\sigma}{N_\sigma},\qquad \|f\|_{2,h}=\left(\Delta x\,\Delta y\sum_{i,j}|f_{ij}|^2\right)^{1/2},\\[0.4em]\delta_{\rm shape}=\frac{\|\widehat n_+-\widehat n_-\|_{2,h}}{\|\widehat n_+\|_{2,h}}=5.55\%.\end{gathered}\]
Норма вычисляется по всем узлам области $256\times256$ мкм без сдвига и поворота. В знаменателе стоит норма нормированного профиля компоненты $+$. На рис.~\ref{fig:3} показаны абсолютные плотности; нормировка используется только для сравнения формы.
% N3_RESULT_END

\reportfigure{N3_spatial_components.png}{Пространственные плотности круговых компонент при кольцевой накачке, \(J=0\) и \(t=2300\) пс. Верхний ряд соответствует \(\varepsilon=0\), нижний — \(\varepsilon=0.3\); левый столбец показывает \(n_+(x,y)\), правый — \(n_-(x,y)\). Все карты приведены на общей абсолютной цветовой шкале в \(\mu\mathrm m^{-2}\). Пунктир обозначает контуры половины максимума накачки. Показана центральная область \(|x|,|y|\leq20\) мкм; населённости \(N_\pm\), подписанные на картах, и поляризация \(Z\), приведённая в тексте, вычислены по всей расчётной области. Параметры кольца: \(r_0=10\) мкм, \(d=2.5\) мкм, \(s_{\mathrm{ring}}=3.16875\). Интегральные населённости и статус установления приводятся в тексте.}{fig:3}



\clearpage
\section[Внутренняя джозефсоновская динамика в 0D-модели]{Внутренняя джозефсоновская динамика\\в 0D-модели}\label{sec:3}

В этой главе рассматривается пространственно однородный (0D) предел уравнения~\eqref{eq:4}. Поля конденсата, плотности резервуаров и скорости накачки не зависят от пространственных координат, поэтому кинетические члены обращаются в нуль. Рисунки~\ref{fig:4}--\ref{fig:6} получены из отдельных однородных расчётов и не описывают пространственную динамику кольцевого конденсата на рис.~\ref{fig:3}.

\subsection{Когерентная связь и обмен частицами}\label{sec:3.1}

В этой главе в уравнении \eqref{eq:4} сохраняется член \(-J\psi_{-\sigma}\). Он описывает когерентную связь двух круговых компонент одного конденсата. При равных диагональных энергиях соответствующее расщепление двух линейно поляризованных состояний составляет \(2J\).

Здесь \(J\) не зависит от волнового вектора. Поэтому для рассматриваемой внутренней джозефсоновской динамики не требуется вводить зависящее от импульса TE–TM-расщепление. Некогерентная накачка задаёт поступление частиц через резервуары; относительная фаза формируется в динамике связанных полей. Внутренняя джозефсоновская связь при нерезонансной накачке рассматривается, в частности, в работе \cite{ref4}.

Две компоненты с круговой поляризацией записываются в виде
\[
\psi_\pm(t)=\sqrt{n_\pm(t)}\,e^{i\phi_\pm(t)}.
\]
Их относительная фаза равна
\[
\theta(t)=\phi_+(t)-\phi_-(t).
\]
Здесь постоянная относительная фаза означает, что \(\theta\) не меняется во времени. Отдельные фазы могут продолжать изменяться, но с одинаковой скоростью. Их разность не обязана быть нулевой. Пространственная однородность предполагается в 0D-модели; она не является результатом расчёта фазовой синхронизации.

Поток между компонентами равен

\begin{equation}\label{eq:9}
I_J
=
\frac{2J}{\hbar}\operatorname{Im}(\psi_-^*\psi_+)
=
\frac{2J}{\hbar}\sqrt{n_+n_-}\sin\theta.
\end{equation}

При таком определении положительный \(I_J\) означает поступление частиц в компоненту \(+\):

\begin{equation}\label{eq:10}
\begin{gathered}
\dot n_+=(Rn_{A,+}-\gamma_C)n_+ + I_J,\\[0.5em]
\dot n_-=(Rn_{A,-}-\gamma_C)n_- - I_J.
\end{gathered}
\end{equation}

Относительная фаза удовлетворяет уравнению

\begin{equation}\label{eq:11}
\dot\theta
=
-\frac{V_+-V_-}{\hbar}
-\frac{2J}{\hbar}
\frac{z}{\sqrt{1-z^2}}\cos\theta.
\end{equation}

Различие энергетических сдвигов вызывает набег относительной фазы. Связь \(J\) одновременно влияет на фазу и перераспределяет частицы, изменяя плотности и сами энергетические сдвиги.

\subsection{Три наблюдаемых режима}\label{sec:3.2}

Расчёты выполняются при \(s=2\), \(\varepsilon=0.3\) и одинаковых начальных условиях. Рассматриваются три значения связи:

\begin{itemize}\setlength{\itemsep}{3pt}
\item \(J=5\) мкэВ: плотности периодически изменяются, относительная фаза продолжает накапливаться;
\item \(J=6.5\) мкэВ: плотности сохраняют колебания, но относительная фаза колеблется в ограниченном диапазоне;
\item \(J=10\) мкэВ: плотности и относительная фаза становятся постоянными во времени.
\end{itemize}

Следовательно, прекращение среднего набега относительной фазы ещё не означает прекращения колебаний плотности.

При \(J=10\) мкэВ стационарный обменный поток составляет

\[
I_J\simeq0.01477\ \mu\mathrm m^{-2}\,\mathrm{ps}^{-1}.
\]

Он компенсирует различие между поступлением и потерями частиц в двух компонентах. Поэтому постоянная относительная фаза допускает ненулевой обменный поток. Сами комплексные поля при этом продолжают вращаться с общей частотой.

\reportfigure{N4_component_dynamics.png}{Пространственно однородные плотности конденсата \(n_+(t)\) и \(n_-(t)\) при трёх значениях когерентной связи: \textbf{(a)} \(J=5\), \textbf{(b)} \(6.5\), \textbf{(c)} \(10\) мкэВ. Во всех случаях \(s=2\), \(\varepsilon=0.3\); начальные условия одинаковы. Показаны поздние интервалы после проверки характера установившегося движения. По горизонтали отложено \(\tau=t-t_0\); во всех панелях \(t_0=16966.5\) пс, конец интервала — \(18000\) пс. Продолжительность показанных интервалов и шкала плотности одинаковы.}{fig:4}





\reportfigure{N5_relative_phase.png}{Относительная фаза \(\theta(t)\) между двумя компонентами с круговой поляризацией в однородной модели для тех же траекторий и временных интервалов, что на рис.~\ref{fig:4}. Панели соответствуют \(J=5,\ 6.5,\ 10\) мкэВ. При слабой связи наблюдается непрерывный набег фазы, при промежуточной — ограниченные фазовые колебания, при сильной — постоянная во времени относительная фаза. Фаза развёрнута непрерывно; для размещения кривых допускается только сдвиг на целое число \(2\pi\), сохраняющий её физическое значение по модулю \(2\pi\).}{fig:5}



\subsection{Переход к постоянным плотностям и относительной фазе}\label{sec:3.3}

Для характеристики зависимости от \(J\) используются две величины:

\begin{equation}\label{eq:12}
A_z=\frac{\max z-\min z}{2},
\qquad
\Delta\overline{\omega}
=\overline{\omega_+-\omega_-}
=-\overline{\dot\theta}.
\end{equation}

Первая величина — амплитуда колебаний дисбаланса. Для второй принимается соглашение \(\omega_\sigma=-\dot{\arg\psi_\sigma}\). При периодическом движении
\[
\Delta\overline{\omega}
=-\left\langle\frac{\theta(t+T)-\theta(t)}{T}\right\rangle_t,
\]
где \(T\) — период, а усреднение проводится внутри проверенного позднего окна. Для стационарных и несвязанных компонент берётся отрицательный наклон относительной фазы.

Средний набег фазы исчезает около \(J=6.4\) мкэВ, но плотности продолжают колебаться. Переход к постоянным плотностям и относительной фазе находится между рассчитанными точками \(6.84\) и \(6.86\) мкэВ.

При \(J=6.38\) и \(6.40\) мкэВ одна компонента на части периода слишком слаба для определения фазы. Периодичность плотностей подтверждена, но \(\Delta\overline{\omega}\) этим точкам не приписывается.

Это смена режима существующего конденсата; её положение зависит от накачки, параметров и подготовки. Все точки запущены независимо из одинакового начального состояния. Полученная зависимость не является полной картой возможных состояний.

\reportfigure{N6_coupling_summary.png}{Смена динамического режима в однородной модели при увеличении связи \(J\); \(s=2\), \(\varepsilon=0.3\). \textbf{(a)} Амплитуда колебаний дисбаланса \(A_z\). \textbf{(b)} Разность средних частот \(\Delta\overline{\omega}\), в \(\mathrm{ps}^{-1}\). Каждая точка — независимый старт с одинаковыми начальными условиями после проверки завершения перехода. Частоты определяются по изменению фазы за полный период с обратным знаком. Серые ромбы при \(J=6.38\) и \(6.40\) мкэВ обозначают циклы с ненадёжной фазой из-за малой плотности компоненты; пропуски частоты не означают нуль. Цветом отмечены \(J=5,\ 6.5,\ 10\) мкэВ с рис.~\ref{fig:4} и~\ref{fig:5}. Стрелкой подписан интервал \(6.84\)--\(6.86\) мкэВ, где прекращается первый плотностный цикл.}{fig:6}



\appendix
\clearpage
\section{Выбор и согласование параметров}\label{sec:A}

\subsection{Параметры прежней модели}\label{sec:A.1}

Базовые коэффициенты сохраняются из предыдущего отчёта, использовавшего модель Alyatkin et al. \cite{ref1}. В частности, \(R\) является параметром моделирования, а соотношение \(g_R=2g_C\) принято в этой модели. Перенос этих значений обеспечивает последовательность сравнения, но не означает, что все коэффициенты независимо измерены для рассматриваемой структуры.

\begin{center}\begin{tabular}{@{}p{0.46\textwidth}p{0.46\textwidth}@{}}\toprule
Величина & Значение\\\midrule
\(\hbar\) & \(0.658212\ \mathrm{meV}\,\mathrm{ps}\)\\
\(m/m_e\) & \(5.63\times10^{-5}\)\\
\(K=\hbar^2/(2m)\) & \(0.676729\ \mathrm{meV}\,\mu\mathrm m^2\)\\
\(g_C=\alpha_1\) & \(0.0024\ \mathrm{meV}\,\mu\mathrm m^2\)\\
\(g_R=g_A\) & \(0.0048\ \mathrm{meV}\,\mu\mathrm m^2\)\\
\(R\) & \(0.021\ \mu\mathrm m^2\,\mathrm{ps}^{-1}\)\\
\(\gamma_C\) & \(0.182\ \mathrm{ps}^{-1}\)\\
\(\gamma_R=\gamma_A\) & \(0.050\ \mathrm{ps}^{-1}\)\\
\(\eta\) & \(9.8\ \mathrm{ps}\)\\
\bottomrule\end{tabular}\end{center}

Все \(n\) имеют размерность \(\mu\mathrm m^{-2}\), а скорости накачки — \(\mu\mathrm m^{-2}\,\mathrm{ps}^{-1}\).

\subsection{Неактивный резервуар}\label{sec:A.2}

Условие согласования \(g_I/W=g_R\eta\) определяет отношение двух новых коэффициентов, но не задаёт их по отдельности. В основном наборе дополнительно принимается \(g_I=g_R\). Тогда

\[
W=\frac{1}{\eta}=0.102041\ \mathrm{ps}^{-1}.
\]

Для собственных потерь неактивного резервуара используется ориентир \(\gamma_I=1/550\ \mathrm{ps}^{-1}\) из модели Gnusov et al. \cite{ref2}. Остальные параметры этой работы целиком не переносятся: \(\gamma_A\) сохраняется из прежнего отчёта, а \(W\) определяется указанным согласованием.

Итоговые значения:

\[
\begin{gathered}
g_I=0.0048\ \mathrm{meV}\,\mu\mathrm m^2,\qquad
\tau_I=\frac{1}{\gamma_I+W}=9.62844\ \mathrm{ps},\\[0.5em]
P_{\mathrm{in}}=1.01782P_A.
\end{gathered}
\]

Время \(550\) пс характеризует потери частиц без учёта перехода в активный резервуар. Отклик неактивного резервуара определяется суммой скоростей потерь и перехода, поэтому составляет около \(9.63\) пс.

\subsection{Спиновые коэффициенты}\label{sec:A.3}

Для взаимодействия поляритонов одинаковой поляризации сохраняется \(\alpha_1=g_C\). Для противоположных поляризаций принимается

\[
\alpha_2=-0.1\alpha_1
=-0.00024\ \mathrm{meV}\,\mu\mathrm m^2.
\]

Отношение \(-0.1\) соответствует модельному выбору в работе Askitopoulos et al. \cite{ref4}. Скорости спинового перераспределения выбираются равными

\[
\Gamma_{s,A}=\Gamma_{s,I}=0.1\ \mathrm{ps}^{-1},
\]

с опорой на масштаб, использованный Gnusov et al. \cite{ref3}. Эти значения служат фиксированными параметрами расширения; они не представлены как результаты измерения нашей структуры.

При записи обменного члена в форме \(\Gamma_s(n_{-\sigma}-n_\sigma)\) разность населённостей из-за одного только обмена затухает со скоростью \(2\Gamma_s\). Соответствующее время равно \(5\) пс.

Система~\eqref{eq:4} точно переходит в скалярную модель с двумя резервуарами~\eqref{eq:2}, если \(J=\Gamma_{s,A}=\Gamma_{s,I}=0\), а накачка, поле конденсата и обе плотности резервуаров компоненты \(-\) равны нулю. Нулевые поля и плотности резервуаров задаются в начальный момент и затем остаются нулевыми. Условия согласования~\eqref{eq:3} сохраняют стационарные решения прежней однорезервуарной модели, но не произвольные переходные процессы. При основных ненулевых \(\Gamma_s\) резервуары перераспределяют частицы между поляризациями.

\subsection{Масштаб когерентной связи}\label{sec:A.4}

Опорное значение \(J=0.010\) мэВ соответствует линейному поляризационному расщеплению \(2J=20\) мкэВ. Расщепления порядка \(15\)–\(20\) мкэВ измерялись Paraïso et al. \cite{ref5} в другой микрорезонаторной структуре. Здесь этот результат используется как ориентир масштаба, а зависимость от \(J\) исследуется отдельно.

Таким образом, новые коэффициенты образуют согласованное феноменологическое расширение прежней модели. Расчёты не заявляются как количественное воспроизведение одного конкретного эксперимента.

\clearpage
\section{Численный метод и проверки}\label{sec:B}

В однородных расчётах пространственные производные опускаются, но сохраняется размерность двумерной плотности. Поля и резервуары интегрируются совместно методом DOP853 с относительным допуском \(2\times10^{-11}\) и абсолютным \(2\times10^{-13}\). Для спиновых серий начальные условия имеют вид

\[
\psi_\pm(0)=\sqrt{5\times10^{-5}}\,e^{\pm0.1i},
\qquad n_{A,\pm}(0)=n_{I,\pm}(0)=0.
\]

Стационарные значения при \(J=0\) проверены по формулам \eqref{eq:6}. Для колебательных режимов проверяются повторяемость плотностей и поляризационного состояния, период и сохранение амплитуды при увеличении времени расчёта. Выбранные точки дополнительно пересчитаны с более строгими параметрами интегрирования.

Для пространственной серии используется периодическая область \(256\times256\) мкм, сетка \(512\times512\) и шаг \(\Delta t=0.025\) пс. Хранятся два комплексных поля и четыре вещественные плотности резервуаров; \(n_\pm\) вычисляются как \(|\psi_\pm|^2\).

Один шаг состоит из кинетического полушага для обоих полей, совместного локального шага RK4 и второго кинетического полушага. Полная схема расщепления имеет второй порядок точности по времени. В пространстве Фурье кинетический полушаг выполняется умножением на

\[
\exp\!\left[
-\frac{iK(k_x^2+k_y^2)\Delta t}{2\hbar}
\right].
\]

Все четыре пространственных расчёта достигли \(2300\) пс. Оба последних окна \(2100\)--\(2200\) и \(2200\)--\(2300\) пс проходят критерии установления: для каждой населённости конденсата и резервуаров размах и дрейф между средними по половинам окна, отнесённые к среднему модулю за окно, менее \(0.5\%\); для каждой карты плотности конденсата и резервуаров относительное L2-изменение от начала окна менее \(1\%\), с проверкой через \(5\) пс. При \(J=0\) условие на относительную фазу не накладывается.

% N3_CHECKS_START
Для $\varepsilon=0.3$ выполнены независимые повторы с половинным временным шагом и грубой пространственной сеткой. Относительные ошибки определяются как
\[\begin{gathered}\delta N_\sigma=\frac{|N_\sigma^{(c)}-N_\sigma^{(\rm ref)}|}{|N_\sigma^{(\rm ref)}|},\\[0.4em]\delta_2(n_\sigma)=\frac{\|n_\sigma^{(c)}-\mathcal S n_\sigma^{(\rm ref)}\|_{2,h_c}}{\|\mathcal S n_\sigma^{(\rm ref)}\|_{2,h_c}}.\end{gathered}\]
Здесь ref --- основной расчёт на $512^2$ с шагом $0.025$ пс, $c$ --- контроль. Оператор $\mathcal S$ выбирает узлы контрольной сетки из эталона и тождественен при сравнении шагов. Норма охватывает всю область на сетке контроля. В таблице приведены $100\delta N_\sigma$ и $100\delta_2(n_\sigma)$ при общем конечном времени, то есть ошибки в процентах.
\begin{center}\small\begin{tabular}{@{}lrrrrc@{}}\toprule
Контроль & $\delta N_+$ & $\delta N_-$ & $\delta_2(n_+)$ & $\delta_2(n_-)$ & Статус\\\midrule
Сетка $256^2$ & $9.67\times10^{-7}$ & $6.03\times10^{-7}$ & $3.27\times10^{-5}$ & $2.95\times10^{-5}$ & PASS\\
Шаг $0.0125$ пс & $1.40\times10^{-3}$ & $1.76\times10^{-3}$ & $1.30\times10^{-3}$ & $1.68\times10^{-3}$ & PASS\\
\bottomrule\end{tabular}\end{center}
Обе ошибки должны быть менее $1\%$ для каждой компоненты. Поворот и подгонка амплитуды не применяются.
Максимальные относительные значения по четырём расчётам: невязка общего баланса $1.18\times10^{-12}$, краевой вес $3.05\times10^{-7}$, высокочастотный вес $2.59\times10^{-10}$. Пороговые значения: $10^{-3}$, $10^{-4}$ и $10^{-6}$ соответственно.
Итоговая оценка новой пространственной серии: \textbf{PASS}.
% N3_CHECKS_END


\clearpage
\begin{thebibliography}{5}
\bibitem{ref1}
S.~Alyatkin et al.
Antiferromagnetic Ising model in a triangular vortex lattice of quantum fluids of light.
\textit{Science Advances} \textbf{10}(34), eadj1589 (2024).
\url{https://doi.org/10.1126/sciadv.adj1589}
\bibitem{ref2}
I.~Gnusov et al.
Quantum vortex formation in the ``rotating bucket'' experiment with polariton condensates.
\textit{Science Advances} \textbf{9}(4), eadd1299 (2023).
\url{https://doi.org/10.1126/sciadv.add1299}
\bibitem{ref3}
I.~Gnusov et al.
Optical orientation, polarization pinning, and depolarization dynamics in optically confined polariton condensates.
\textit{Physical Review B} \textbf{102}, 125419 (2020).
\url{https://doi.org/10.1103/PhysRevB.102.125419}
\bibitem{ref4}
A.~Askitopoulos et al.
Nonresonant optical control of a spinor polariton condensate.
\textit{Physical Review B} \textbf{93}, 205307 (2016).
\url{https://doi.org/10.1103/PhysRevB.93.205307}
\bibitem{ref5}
T.~K.~Para\"iso et al.
Multistability of a coherent spin ensemble in a semiconductor microcavity.
\textit{Nature Materials} \textbf{9}, 655--660 (2010).
\url{https://doi.org/10.1038/nmat2787}
\end{thebibliography}
\end{document}
